\section{机制解释的对照}\label{sec:mechanism}

信号层与统计层的比较完成之后，还有一个问题悬着：\emph{为什么会这样}。本章对照两条脉络对机制的争夺。这不是学理上的装饰——机制决定了策略在什么条件下应当被信任、在什么条件下应当被放弃。

\subsection{两套解释框架}

文献中对收益可预测性的解释可以粗分为两族。

\textbf{风险补偿族}认为，可预测性来自时变的风险溢价：某些状态下的预期收益更高，是因为那些状态下风险更大。动量在危机中获利，说明它承担了被市场厌恶的风险；反转策略的收益来自承担流动性风险与资金约束风险。

\textbf{行为偏差族}认为，可预测性来自投资者的系统性认知偏差：反应不足产生中短期动量，反应过度产生长期反转。

两族解释在文献中长期并存。Cochrane(2011) 的贴现率视角提供了一个把二者装进同一框架的努力：收益率可预测性的核心是贴现率的时变性，而贴现率的变化既可以是风险溢价的理性变化，也可以是情绪驱动的错价\cite{cochrane2011}。Cochrane(2005) 的综述进一步指出，区分二者的困难在于\emph{两者都能拟合同样的预测回归结果}\cite{cochrane2005}。

\subsection{反转一侧的机制}

反转策略的机制解释有三个层次。

\textbf{第一层是过度反应}。De Bondt \& Thaler(1985) 的原始解释：投资者对坏消息过度悲观，把价格压到低于合理水平，随后价格回升\cite{debondt1985}。Summers(1986) 与 Shiller(1981) 的价格波动过大证据为这一解释提供了宏观背景\cite{summers1986,shiller1981}。Lakonishok 等(1994) 把这一逻辑推广到价值投资：反向投资之所以有效，是因为市场外推了过去的增长\cite{lakonishok1994}。

\textbf{第二层是噪声交易者风险}。De Long 等(1990) 的模型给出了一个更精细的机制：即便理性套利者知道价格偏离，他们也不会立即纠正，因为噪声交易者可能让偏离进一步扩大，套利者承担了额外的风险\cite{delong1990}。这一模型解释了\emph{为什么偏离能持续}——而这正是反转策略需要“等待”的原因。

\textbf{第三层是流动性补偿}。这是最被后来的实证支持的一层。Lo \& MacKinlay(1990) 指出大量“反转利润”实为交叉自相关造成的聚合假象\cite{lomackinlay1990}；短周期的反转收益与买卖价差、流动性提供高度绑定。Avellaneda \& Lee(2010) 的因子剥离框架把反转定义为“残差的均值回复”，其实质是在剥离共同因子暴露后交易特质性错价\cite{avellaneda2010}。流动性作为定价因子的地位由 Pástor 与 Stambaugh(2003) 系统化：小盘股对市场流动性冲击更敏感，流动性风险的补偿构成收益的一部分\cite{pastorstambaugh2003}。这一层也解释了反转策略为什么容易在资金端出问题——Shleifer 与 Vishny(1997) 之后，Gromb 与 Vayanos(2010) 对该族理论做了系统综述\cite{shleifervishny1997,grombvayanos2010}。

三层机制并不互斥，而是作用在不同时间尺度上：短尺度以流动性补偿为主，中尺度混合过度反应与噪声交易者风险，长尺度以估值回归为主。这恰好对应第~\ref{sec:duality}~节“尺度决定归属”的判断。

\subsection{趋势一侧的机制}

趋势策略的机制解释同样分层。

\textbf{延迟反应}。Hong \& Stein(1999) 的信息渐进扩散模型是这一族最完整的表述：市场由“消息观察者”与“动量交易者”两类参与者构成，前者缓慢吸收私有信息，后者仅从价格推断信息，二者的互动先产生反应不足、后产生过度反应\cite{hongstein1999}。Barberis 等(1998) 用保守主义与代表性启发两个偏差给出平行解释\cite{bsv1998}；Daniel 等(1998) 用过度自信与自我归因给出第三个版本\cite{dhs1998}。三篇文献的构造不同，但都产出“先动量后反转”的同一预测——这与第~\ref{sec:unity}~节的自相关 U 形结构吻合。

\textbf{风险溢价}。Hurst 等(2017) 的百年证据是这一解释的主要支撑：趋势跟踪在每个十年、各种宏观环境下都为正，且与股债的长期相关性低\cite{hop2017}。若趋势收益纯属错价，它应当随市场成熟而消失；百年尺度上的稳定性更符合“承担了某种系统性风险”的解释。Erb \& Harvey(2006) 从商品展期收益的角度提供了一个具体的风险来源候选\cite{erbharvey2006}。Ilmanen(2011) 对各类风险溢价的系统梳理为这一族解释提供了统一框架：动量的收益可以按“承担了何种被厌恶的风险、在何种环境下兑现”来归类\cite{ilmanen2011}。一个有助于区分机制的线索是：若动量只是“对近期信息反应不足”的表象，那么它与基本面动量的相关性应当很高；Novy-Marx(2012) 检验后发现并非如此，动量并非主要由基本面动量驱动\cite{novymarx2012}。

\textbf{拥挤与协调失败}。Stein(2009) 的论证是第三种解释的雏形：动量属于“未锚定策略”（unhedged strategy），套利者无法知道同行的建仓规模，因而无法判断价格偏离是否会被纠正\cite{stein2009}。这一机制不解释趋势为何长期有效，但解释它为何会突然失效——这与动量崩溃的实证相符\cite{danielmoskowitz2016}。

\subsection{一个机制生成两个符号}

本章最重要的观察是：\textbf{不需要两套机制来解释两策略}。行为模型（Hong--Stein、BSV、DHS）在单一框架内同时产生短期动量的正自相关与长期反转的负自相关。这意味着第~\ref{sec:genealogy}~节观察到的“分岔”在机制层面根本不存在——两条脉络是在\emph{同一个机制的产物}上各自取了不同的一端。

这一观察有直接的操作含义。如果两策略来自同一机制，那么它们的失效也应当同时发生，因为机制被破坏时两个方向都会失效。经验证据支持这一点：2007 年 8 月，量化反转基金与趋势基金同时遭受冲击，因为冲击来自共同的资金端而非各自的信号端\cite{khandanilo2011,brunnermeier2009}。这解释了为什么简单地把两个策略等权组合起来并不能提供想象中的分散化——它们在某些状态下的相关性会趋近于一。

\subsection{机制之争对策略的实际意义}

机制解释的分歧不是纯粹的学术趣味，它决定了三件事。

\textbf{一是对“样本外失效”的解释}。若趋势收益是风险溢价，发表后不应衰减（风险仍需承担）；若是错价，则应当随套利而衰减。McLean \& Pontiff(2016) 发现的发表后衰减支持后者\cite{mcleanpontiff2016}；Hurst 等(2017) 的百年稳定性支持前者\cite{hop2017}。两种证据并存，说明趋势收益可能同时含两个成分，且比例随时间变化。

\textbf{二是对“何时应该收手”的判断}。若反转收益主要是流动性补偿，那么在流动性环境改善（如算法做市普及）后应当衰减——这与 Do \& Faff(2012) 报告的配对交易利润衰减时序吻合\cite{dofaff2012}。

\textbf{三是对“该不该加杠杆”的判断}。若趋势收益是风险溢价，加杠杆应当获得成比例的更高收益（在无摩擦条件下）；若是错价，加杠杆会放大拥挤风险。Stein(2009) 的杠杆外溢论证正是针对后者的警告\cite{stein2009}。

\subsection{边界：机制不可直接观测}

必须承认本章讨论的机制都无法被直接观测。我们观测到的只是收益序列与自相关结构，机制是\emph{对这些观测的解释}。Cochrane(2011) 的说法值得引述：几乎所有资产定价结果都可以同时用理性风险与行为偏差来解释，因为两者在价格数据上常难以区分\cite{cochrane2011}。

因此本章的结论应当保守表述：\emph{关于机制的争论应当影响我们对策略稳健性的先验，而不应当成为选择策略的依据}。选择依据应当是可以检验的东西——自相关结构在哪一段尺度上显著、成本与容量的实测约束、以及资金端的状态。这也是本文把重心放在第~\ref{sec:execution}~节与第~\ref{sec:empirical}~节的原因。

\subsection{能否设计一个可判决的检验}

机制不可直接观测，不等于无法设计区分性的检验。可行路径有三条，各自的可行性与困难值得写清，因为它们构成第~\ref{sec:future}~节遗留课题 Q9 的具体内容。

\textbf{路径一：看样本外是否衰减}。若趋势收益是风险溢价，它在文献发表后\emph{不应}系统性衰减（风险仍需有人承担，补偿仍在）；若是错价，则应当随套利而衰减。McLean 与 Pontiff(2016) 的实测支持后者\cite{mcleanpontiff2016}；Hurst 等(2017) 的百年稳定性支持前者\cite{hop2017}。困难在于两者可以同时成立——收益中同时含溢价成分与错价成分，比例还随时间变化，这使得“衰减多少”无法直接读成“溢价占多少”。

\textbf{路径二：看辅助变量的预测方向}。若反转收益来自流动性补偿，那么它与流动性指标（买卖价差、成交量、资金成本）应当同向变动；若来自过度反应，则应与情绪指标同向。Pástor 与 Stambaugh(2003) 的流动性因子框架提供了可用的代理\cite{pastorstambaugh2003}，Brunnermeier 与 Pedersen(2009) 的融资流动性模型给出了理论预测方向\cite{brunnermeier2009}。这条路径的困难在于代理变量的内生性：流动性差本身就是价格偏离的一部分。

\textbf{路径三：用状态的不可预测性做自然实验}。若收益是风险补偿，它在事前可识别的“高风险状态”下应当更高；若是错价，其高低应与是否被套利相关，而与事前风险状态关系较弱。Daniel 与 Moskowitz(2016) 对恐慌状态的刻画、以及 Ang 与 Timmermann(2012) 综述的 regime 方法提供了状态划分工具\cite{danielmoskowitz2016,angtimmermann2012}。困难是状态划分本身依赖模型，而不同模型给出的状态划分会导致不同的检验结论——这构成了近乎循环的依赖。

三条路径共同的障碍是同一个：\emph{任何用于区分机制的变量，都同时被机制所影响}。这解释了为什么这一问题在几十年后仍未闭合。本文的态度是把它记录为开放的困难，而不给出一个虚假的裁决。
